\hypertarget{ej8}{}
\noindent \textbf{Ejercicio 8.} Determinar, para cada aparición de variables, si dicha aparición se encuentra libre o ligada. En caso de estar ligada, aclarar a qué cuantificador lo está. En los casos en que sea posible, proponer valores para las variables libres de modo tal que las expresiones sean verdaderas.

\begin{enumerate}[label=\alph*)]
    \item $(\forall x: \mathbb{Z}) (0 \leq x < n \rightarrow x + y = z)$
    \item $(\exists x: \mathbb{Z}) ((\forall y: \mathbb{Z}) (0 \leq x < n \wedge 0 < y < m \rightarrow x + y = z))$
    \item $(\forall j: \mathbb{Z}) (0 \leq j < 10 \rightarrow j < 0)$
    \item $(\forall j: \mathbb{Z}) (j \leq 0 \rightarrow P(j)) \wedge P(j)$
\end{enumerate}

\vspace{0.2cm}
{\color{gray}\hrule}
\vspace{0.4cm}

\textbf{Resolución:}

Recordemos la teoría básica: una variable está \textbf{ligada} si se encuentra bajo el alcance de un cuantificador (como $\forall$ o $\exists$) que la referencia explícitamente. Está \textbf{libre} si su valor depende del contexto exterior y no está atada a ningún cuantificador en la expresión.

\begin{enumerate}[label=\textbf{\alph*)}]

    \item \textbf{$(\forall x: \mathbb{Z}) (0 \leq x < n \rightarrow x + y = z)$} \\
    \textbf{Análisis de variables:}
    \begin{itemize}
        \item $x$: \textbf{Ligada} al cuantificador universal $(\forall x: \mathbb{Z})$.
        \item $n$, $y$, $z$: \textbf{Libres}.
    \end{itemize}
    \textbf{Valores para que sea verdadera:} Proponemos $n = 0$, con $y, z \in \mathbb{Z}$ arbitrarios. Al ser $n=0$, la condición del antecedente $0 \leq x < 0$ siempre es falsa. Como un antecedente falso hace que la implicación completa sea Verdadera (vacuidad), la expresión es Verdadera.

    \vspace{0.4cm}
    \item \textbf{$(\exists x: \mathbb{Z}) ((\forall y: \mathbb{Z}) (0 \leq x < n \wedge 0 < y < m \rightarrow x + y = z))$} \\
    \textbf{Análisis de variables:}
    \begin{itemize}
        \item $x$: \textbf{Ligada} al cuantificador existencial $(\exists x: \mathbb{Z})$.
        \item $y$: \textbf{Ligada} al cuantificador universal $(\forall y: \mathbb{Z})$.
        \item $n$, $m$, $z$: \textbf{Libres}.
    \end{itemize}
    \textbf{Valores para que sea verdadera:} Usamos la misma lógica del inciso anterior. Proponemos $m = 0$ (con $n, z \in \mathbb{Z}$ arbitrarios). De esta forma, el antecedente $0 < y < 0$ es absurdo (Falso), la implicación se vuelve trivialmente verdadera, y por lo tanto existe cualquier $x$ que la cumple.

    \vspace{0.4cm}
    \item \textbf{$(\forall j: \mathbb{Z}) (0 \leq j < 10 \rightarrow j < 0)$} \\
    \textbf{Análisis de variables:}
    \begin{itemize}
        \item $j$: \textbf{Ligada} al cuantificador universal $(\forall j: \mathbb{Z})$.
    \end{itemize}
    \textbf{Valores para que sea verdadera:} Esta expresión \textbf{no tiene variables libres} a las que asignarles valores (es una fórmula cerrada). Además, evalúa siempre a \textbf{Falsa}. (Por ejemplo, si instanciamos $j=1$, el antecedente $0 \leq 1 < 10$ es $V$, pero el consecuente $1 < 0$ es $F$, resultando en $V \rightarrow F \equiv F$). Por lo tanto, no es posible proponer valores.

    \vspace{0.4cm}
    \item \textbf{$(\forall j: \mathbb{Z}) (j \leq 0 \rightarrow P(j)) \wedge P(j)$} \\
    \textbf{Análisis de variables:} \\
    ¡Atención acá con el alcance de los paréntesis!
    \begin{itemize}
        \item $j$ (primeras dos apariciones): \textbf{Ligada} al cuantificador $(\forall j: \mathbb{Z})$. El alcance termina al cerrar el paréntesis de la implicación.
        \item $j$ (última aparición en el último $P(j)$): \textbf{Libre}. Queda afuera del cuantificador.
        \item $P$: \textbf{Libre} (es un símbolo de predicado libre que debemos instanciar).
    \end{itemize}
    \textbf{Valores para que sea verdadera:} Proponemos que el predicado libre sea una tautología, por ejemplo $P(x) = True$ para todo $x$, y que la variable libre sea $j = 1$. Así, la implicación $(\forall j) (j \leq 0 \rightarrow True)$ es verdadera, y el término suelto $P(1)$ también es verdadero. Resultando en $True \wedge True \equiv True$.

\end{enumerate}

\vspace{0.5cm}
\hrule
\vspace{0.5cm}